\section{Key Technical Results: Numerical Guarantees for Double-Word
Operations over IEEE~754}

\subsection{From Idealized Arithmetic to IEEE~754}
\label{sec:flx-to-ieee}

The existing mechanized analyses of double-word arithmetic are stated in
Flocq's FLX format \citep{Rideau2022}.  FLX provides fixed precision over an
unbounded exponent range and therefore models neither gradual underflow nor
overflow.  We use Flocq's FLT format to model gradual underflow and establish
separate upper-range bounds to prove that the absence of overflow.

Flocq defines a floating-point format by the real numbers it represents.  An
integer significand $m$ and integer exponent $q$ represent the real value
$m\beta^q$.  For radix $\beta>1$ and precision $p>0$, FLX and FLT are defined
as follows.
\begingroup
\small
\begin{flalign}
  &\mathbb{F}_{\mathrm{FLX}}(\beta,p)
    =
    \left\{m\beta^q \mathrel{}\middle|\mathrel{}
      m,q\in\mathbb{Z},\ |m|<\beta^p\right\}
    &&
    \label{eq:flx-format}\\
  &\mathbb{F}_{\mathrm{FLT}}(\beta,p,e_{\min})
    =
    \left\{m\beta^q \mathrel{}\middle|\mathrel{}
      m,q\in\mathbb{Z},\ |m|<\beta^p,\ e_{\min}\leq q\right\}
    &&
    \label{eq:flt-format}
\end{flalign}
\endgroup
The lower bound $e_{\min}\leq q$ is the only additional constraint in FLT
\citep[Sections~III-C and III-D]{Boldo2011}.  Even when the same rounding
rule is applied in both formats, the rounded results may differ because the
representable values have different spacing near zero.  Here, spacing refers
to the distance between adjacent representable values at a given magnitude.

In FLX, the spacing continues to decrease as the magnitude decreases.  FLT
has the same spacing as FLX down to its smallest positive normal value
$\nu=\beta^{e_{\min}+p-1}$.  Below $\nu$, the exponent cannot decrease below
$e_{\min}$, so the representable values have a fixed spacing of
$\eta=\beta^{e_{\min}}$.  The value $\eta$ is also the smallest positive
value representable in FLT.  This fixed spacing enables gradual underflow:
values below $\nu$ remain representable as multiples of $\eta$, but with fewer
significant bits as they approach zero.  The difference in spacing may cause
FLT and FLX rounding to produce different results.
For binary64, $\eta=2^{-1074}$ and $\nu=2^{-1022}$.

For round-to-nearest, ties-to-even, rounding a value whose magnitude is at
least $\nu$ gives the same result in both formats.  Let
$\operatorname{RN}_{\mathrm{FLX}}$ and
$\operatorname{RN}_{\mathrm{FLT}}$ denote rounding in FLX and FLT,
respectively.  Then for an exact real value $x$
\begin{equation}
  |x|\geq\nu
  \quad\Longrightarrow\quad
  \operatorname{RN}_{\mathrm{FLT}}(x)
    =\operatorname{RN}_{\mathrm{FLX}}(x)
  \label{eq:round-flt-flx}
\end{equation}
For double-word arithmetic in the normal range, this equality lets us reuse
the analyses of \citet{Joldes2017} as mechanized in FLX by
\citet{Rideau2022}.  For
$0<|x|<\nu$, \cite{eq:round-flt-flx} does not apply.  However, the
 roundings may nevertheless be equal, so
we either prove their equality separately or reason directly about FLT
rounding.

The FLT format accounts for gradual underflow but it does not account
overflow.  Neither FLX nor FLT has a largest representable real because the
exponent $q$ remains unbounded above.  Binary64 instead has a largest finite
value and may produce an infinity when an operation exceeds its finite range
\citep{IEEE754-2019}.  We therefore prove operation-specific upper-range
bounds that keep each relevant binary64 result finite.  Under these bounds,
the correctness result for the binary64 operation identifies the real value
of its result with the corresponding FLT rounding.

Relating an FLX proof to IEEE~754 therefore creates two obligations at each
rounded expression.  We must prove that the corresponding operation produces
a finite result in the target format, and we must determine whether FLT and
FLX rounding agree.  When they do not, the required numerical property must
be proved directly in FLT.
